% Modernised reading — fonds Alexandre Grothendieck, shelfmark 27, pages 1-36.
%
% This is an INTERPRETATION, not a transcription. Everything the critical
% apparatus carried moves into footnotes. In any dispute of substance the
% transcription governs: whoever cites Grothendieck must cite batch-01.fr.tex
% and batch-02.fr.tex.
%
% Pass: Opus 5.5 (claude-opus-5-5), 2026-09-24 — first pass, unchecked against
% the pages by a human. The folder was read whole: two batches, pages 1-20 and
% 21-36, both transcribed. Pages 1 and 27 are covers (their inscriptions become
% headings), 7 and 16 are blank, 20 and 30 are typed versos with nothing of
% his; nothing here reads them beyond a mention. Opus 5.5 is held to the same
% standard as the Opus 5 reference reading of folder 115 but was never measured
% against it.
%
% What the folder is. A chemise « Divers pour XII à XIV » of complements to
% SGA 3, in five runs: (1) pages 2-6, five annotated leaves of the typescript
% of Exposé XIII (Prop. 1.3, Cor. 1.4-1.6, Rem. 1.7: G acting on V, W subscheme,
% psi : X = G x^N W -> V, W-regular points), summarised and not re-set, as in
% the transcription; (2) pages 8-10, « Complément XII »: plan of a §6 on the
% maximal central torus, Rosenlicht's unit lemma and its corollaries, the
% universal quotient of multiplicative type, extensions by a torus; (3) pages
% 11-12: extension of subgroups of multiplicative type (XII 5.8-5.9, XI 4.?);
% (4) pages 13-15: Picard groups of affine algebraic groups; (5) pages 17-19
% and 21-26: the typescript's setting specialised to the adjoint action
% (V = W(g), W = W(h)) — eleven conditions on a subalgebra h, a lemma on
% nilpotent Lie algebras, counterexamples in characteristic 2 and over F_p, a
% non-representable centre; (6) pages 27-36, « Théorème de semi-continuité »
% for the numerical invariants of the fibres of a smooth group scheme.
%
% Notation fixed for the whole document. G_m, G_a; X^*(T) for characters (the
% page's bold D(T)). In the typescript and the Lie-algebra run: N the
% stabiliser of W, M_a the transporter of a into W, n = Lie N, m_a = Lie M_a,
% u the normaliser of h in g, W(g) the vector group of g. The universal
% multiplicative quotient of pages 8-10 is written G^tm (page: H). In the
% semicontinuity run the pages change notation between pages 28, 29 and 31;
% the reading fixes p.29's: rho_ab, rho_r, rho_rig = rho_ab + rho_r, rho_ss,
% d_ss, d_rad, d_aff (p.28's and p.31's rho_aff), rho_u, rho_na (p.29's
% rho_naff), rho_n, d; G_aff for the affine part (p.28's H); unprimed = special
% fibre, primed = generic. « Préschéma » is written « schéma ».
%
% Substantive departures from the pages, each footnoted where it occurs:
% — p.4: the gothic letters read m = n_a; read n = m_a as on p.2.
% — p.6 (as summarised in the transcription): « a W-regular if phi^{-1}(a) is
%   finite »; phi^{-1}(a) ≅ M_a contains a copy of N, so this is psi^{-1}(a).
% — p.8: the normalisation of Rosenlicht's u is overwritten; read u(a,b) = 1.
% — p.10: « et lisse », added in pencil, is necessary (mu_p); the N.B. on
%   imperfect fields is right, and the reading gives an example (R_{k'/k} G_m).
% — p.12: the relations between H(n), H(m), H are illegible; reconstructed as
%   _m H(n) = H(m), _n H = H(n) from pp.33-34, where the same family appears.
% — p.14, 1°): as read it would give Pic(PGL_2) = 0 over an algebraically
%   closed field; the reading states what [k':k] actually kills.
% — p.14, 2°): Pic(G) ≅ Pic(G_{k~}) needs Galois invariants (outer form of
%   PGL_3); the reading states Pic(G) ≅ Pic(G_{k_s})^Gamma.
% — p.21: psi birational for SL_2 in characteristic 2 is right, but only
%   because N is not smooth; the reading says why.
% — p.33: the overbrace Z over C/Cent(C)°_red; with Z the quotient the second
%   assertion is false; read Z = Cent(C)°_red.
% — p.35: pi_1(s_0) ≅ Z read as Z^ (profinite).
% — p.36: « un réductif se spécialise en un réductif » is false as it stands.
% — p.36: the « n-rang » bound holds for the whole ell-adic tower, not for
%   one n ((Z/2)^2 in PGL_2).
% The reconstruction of the Frobenius-weight argument (pp.35-36) and of the
% proofs on p.10 and p.18 are ours and footnoted as such.
%
% Modern names given and footnoted as ours: lemme de Rosenlicht (on the page),
% quotient de type multiplicatif maximal, restriction de Weil, lemme de Sansuc,
% sous-algèbre de Cartan, élément régulier, théorème d'Engel, module de Tate,
% poids de Frobenius (Weil), variété semi-abélienne, SGA 7 IX.
%
% What the folder announces and never establishes: §6.1-6.3 (a plan); the
% cancelled 6.4; the Hochschild-cochain sequence of p.10; the cancelled
% Cor. 5.10; the proof of XI Cor. 4.?; conditions c)-d) of p.14; the list of
% fourteen conditions of pp.17 and 19; the boxed remark of p.19; the
% counterexample of p.26; the N.B. on representability (p.25); the point
% « à vérifier » of p.33; the equivalence (3) <=> (9); the reduction to a base
% of finite type over Z (p.35).
\documentclass[11pt,a4paper]{article}
% Le préambule commun à toutes les transcriptions du fonds Grothendieck.
%
% Il tient en une idée : **l'incertitude se compose, elle ne se résout pas.**
% Une transcription qui lisse un mot douteux a détruit la seule chose pour
% laquelle on la faisait — un lecteur doit pouvoir savoir, à chaque endroit,
% ce qui était sur le feuillet et ce qui a été ajouté. Les six macros qui
% suivent sont donc l'appareil critique, et non de la décoration.
%
% Les mêmes commandes sont reconnues par `scripts/render.mjs`, qui produit la
% vue de lecture du site. Une macro ajoutée ici doit l'être aussi là-bas,
% faute de quoi le rendu échouera bruyamment — ce qui est voulu : un rendu
% silencieusement faux est pire qu'un rendu absent.

\ProvidesPackage{grothendieck}[2026/08/08 Transcriptions du fonds Grothendieck]

\usepackage{amsmath,amssymb,amsthm}
\usepackage{mathrsfs}
\usepackage{stmaryrd}
% Les diagrammes commutatifs sont partout dans ce fonds : c'est la langue
% même de la géométrie algébrique de Grothendieck, et une transcription qui
% les rendrait en prose ne transcrirait rien.
\usepackage{tikz-cd}
% Français par défaut — c'est la langue des feuillets — mais l'anglais est
% chargé aussi : la traduction et le résumé sont des documents anglais, et la
% typographie française (espaces avant les deux-points) y serait une faute.
% Ces documents basculent par \selectlanguage{english} en tête de corps.
\usepackage[english,french]{babel}
\usepackage{fontspec}
\usepackage{geometry}
\geometry{a4paper,margin=25mm,marginparwidth=28mm}
\usepackage{marginnote}
\usepackage{xcolor}
% ulem, not soul. `soul`'s \st and \ul are built on a character-by-character
% scan that XeLaTeX's font machinery breaks: the document fails with an
% undefined control sequence rather than degrading. ulem does the same two jobs
% and is Unicode-engine safe. `normalem` stops it hijacking \emph, which the
% transcriptions use for Grothendieck's own emphasis.
\usepackage[normalem]{ulem}
\usepackage{hyperref}
\hypersetup{colorlinks=true,linkcolor=black,urlcolor=black,pdfborder={0 0 0}}

% --- Métadonnées du lot -----------------------------------------------------
% Reprises telles quelles par le script de rendu, qui les affiche en tête de
% la vue de lecture. Elles fixent ce que le fichier prétend couvrir : sans
% elles, une transcription flotte sans point d'attache dans un fonds de seize
% mille feuillets.

\newcommand{\folder}[1]{\gdef\grothFolder{#1}}
\newcommand{\batch}[1]{\gdef\grothBatch{#1}}
\newcommand{\leaves}[2]{\gdef\grothFirst{#1}\gdef\grothLast{#2}}
\newcommand{\foldertitle}[1]{\gdef\grothFoldertitle{#1}}
\gdef\grothFolder{?}\gdef\grothBatch{?}\gdef\grothFirst{?}\gdef\grothLast{?}\gdef\grothFoldertitle{}

% --- Le résumé --------------------------------------------------------------
%
% Il ouvre la lecture modernisée, sous le titre « Résumé », et il est composé
% en petit. Ce n'est pas un ornement : c'est le seul passage du document qui
% ne soit pas des mathématiques, et un lecteur doit pouvoir le reconnaître
% avant de l'avoir lu — puis le sauter s'il connaît déjà le sujet, ou s'y
% arrêter s'il ne le connaît pas. Le filet qui le suit marque la reprise du
% corps.

\newenvironment{resume}
  {\small\setlength{\parskip}{0.45em}}
  {\par\medskip\hrule\medskip}

% --- Mots-clés ---------------------------------------------------------------
%
% En anglais, en clôture du résumé de la lecture modernisée : le vocabulaire
% moderne sous lequel chercher ce que les feuillets construisent. Le manifeste
% du site (`npm run manifest`) les extrait pour étiqueter la cote entière —
% cette ligne est la source unique des tags, il n'y a pas de fichier à côté.
\newcommand{\keywords}[1]{\par\smallskip\noindent
  {\footnotesize\textsc{Keywords} --- \textit{#1}}\par}

% --- L'appareil critique ----------------------------------------------------

% Le numéro de feuillet, en marge. C'est l'ancre qui permet de régler un
% désaccord sur une formule en regardant *un* feuillet plutôt qu'une chemise
% de six cents. Le site s'en sert aussi pour tourner les pages du fac-similé
% quand on fait défiler la transcription.
\newcommand{\leaf}[1]{%
  \marginnote{\footnotesize\color{gray!70}\textsf{#1}}%
  \label{leaf:#1}\ignorespaces}

% Illisible : on ne devine pas. Un mot inventé qui se lit comme les autres est
% le pire résultat possible.
\newcommand{\ill}{\textcolor{red!60!black}{[\,\ldots\,]}}

% Lecture douteuse : le mot est donné, et signalé comme incertain.
\newcommand{\uncertain}[1]{\uline{#1}}

% Ajout de l'éditeur — une lettre manquante, un mot sous-entendu.
\newcommand{\add}[1]{\textcolor{blue!50!black}{[#1]}}

% Biffé par Grothendieck lui-même. Ce qu'il a rayé fait partie du document :
% une rature dit souvent où la pensée a changé de direction.
\newcommand{\struck}[1]{\sout{#1}}

% Note du transcripteur, sur l'état matériel du feuillet.
\newcommand{\note}[1]{\par\smallskip{\small\color{gray!60!black}\textit{[#1]}}\par\smallskip}

% Note marginale de Grothendieck, distincte du corps du texte.
\newcommand{\marginal}[1]{\par\smallskip\noindent\hspace{1em}%
  {\small\color{gray!60!black}#1}\par\smallskip}

% --- Titre ------------------------------------------------------------------

\AtBeginDocument{%
  \begin{center}
    {\large\bfseries Fonds Alexandre Grothendieck --- cote n\textsuperscript{o}~\grothFolder}\\[2pt]
    {\itshape\grothFoldertitle}\\[4pt]
    {\small feuillets \grothFirst--\grothLast\ (lot \grothBatch)}\\[2pt]
    {\footnotesize\color{gray!60!black}Transcription établie d'après le fac-similé numérisé
     par l'Université de Montpellier.\\
     Les lectures incertaines sont soulignées, les illisibles notées [\,\ldots\,],
     les ajouts entre crochets.}
  \end{center}
  \vspace{1em}}

\endinput


\folder{27}
\pages{1}{36}
\dating{[vers 1962-1970]}
\watermark{Édition de démonstration\\interprétation personnelle de l'œuvre}
\foldertitle{SGA D [SGA 3]. Compléments XII à XVII : tapuscrit annoté (s.d.), notes manuscrites (s.d.) — lecture modernisée du dossier entier}

\begin{document}

\section*{Résumé}

\begin{resume}
Ces trente-six feuillets sont des chutes d'atelier : ce qui reste sur la table
quand on a fini d'écrire un gros traité, ici le séminaire sur les
\emph{schémas en groupes}, rédigé au début des années soixante sous le sigle
SGA~3. La chemise porte « Divers pour XII à XIV » : des compléments aux
exposés de ce traité qui étudient les tores, les sous-groupes de Cartan et les
éléments réguliers. On y trouve quelques feuillets dactylographiés qu'il
corrige à la plume, et surtout des notes rapides, à l'encre, où il éprouve des
énoncés.

Un groupe algébrique est un groupe dont les éléments forment une variété, les
opérations étant données par des polynômes : les matrices inversibles, le
cercle, une courbe elliptique. La géométrie y distingue trois matières. Il y a
la matière \emph{compacte}, celle des variétés abéliennes, dont la courbe
elliptique --- une bouée, un tore au sens des topologues --- est le premier
exemple. Il y a la matière \emph{circulaire}, celle des groupes
multiplicatifs : le groupe des nombres complexes non nuls, qui se rétracte sur
un cercle, et ses puissances, qu'on appelle aussi, pour compliquer, des tores.
Et il y a la matière \emph{rectiligne}, celle des groupes unipotents : la
droite munie de l'addition. Tout groupe algébrique se fabrique à partir de ces
trois matières, et plusieurs nombres mesurent combien il en contient.

Le dernier tiers du dossier pose une question sur ces nombres. Qu'arrive-t-il
quand un groupe \emph{dégénère} --- quand il varie dans une famille et qu'en
un point de la famille il change brusquement de forme ? Une courbe elliptique
peut dégénérer en le groupe multiplicatif : on pince un des lacets de la
bouée, et il reste un cylindre. Le groupe multiplicatif peut à son tour
dégénérer en la droite additive : le cercle s'ouvre. Le « théorème de
semi-continuité » des pages 27 à 36 dit que la dégénérescence va toujours dans
ce sens-là et jamais dans l'autre : en se spécialisant, la matière compacte ne
peut que diminuer, la somme des matières compacte et circulaire aussi, la
matière rectiligne ne peut que croître. La démonstration est belle. Pour
montrer qu'une partie compacte ne peut pas surgir, elle compte les points
d'ordre fini, qui se transportent d'un membre de la famille à l'autre, et
regarde comment les symétries d'un corps fini les permutent. Sur le cercle,
ces symétries agissent comme sur les racines de l'unité. Sur une variété
abélienne, elles agissent autrement, et un théorème célèbre de Weil sur les
courbes et les variétés abéliennes des corps finis dit exactement en quoi.
Les deux comportements ne peuvent pas se confondre, et cela suffit.

La partie centrale du dossier traite d'une autre question, plus ancienne :
comment reconnaître, dans un groupe, les éléments « génériques », dont un
voisinage est aussi simple que possible ? Pour les matrices, ce sont celles
dont les valeurs propres sont distinctes. Le traité avait donné une réponse en
toute caractéristique, et ces notes cherchent jusqu'où elle tient. Elles
dressent une liste de onze conditions qu'on pourrait vouloir prendre pour
définition, et les séparent l'une de l'autre par des contre-exemples. Ceux-ci
se trouvent en caractéristique 2, où $1 + 1 = 0$ et où les matrices de trace
nulle ne se comportent plus comme d'habitude, ou sur un corps fini, qui peut
être trop petit pour contenir un seul élément générique. Sur le corps à $p$
éléments, le plan n'a que $p+1$ droites passant par l'origine, et chacun de
ses points est sur l'une d'elles : aucun point n'évite toutes les droites
interdites. Le dossier contient aussi des énoncés plus courts. Le lemme de
Rosenlicht dit qu'une fonction qui ne s'annule jamais sur un produit se
factorise en un produit de deux fonctions. Un calcul du groupe de Picard des
groupes algébriques mesure, par exemple, pourquoi le groupe des matrices à
homothétie près n'est pas simplement connexe. Enfin, un groupe sur un anneau
de valuation discrète a un centre qu'aucun schéma ne représente.

Ce que ces feuillets montrent peut-être le mieux, c'est une manière de
travailler. Un énoncé est écrit, puis son contre-exemple est cherché dans les
plus petits groupes qui soient, $\mathrm{SL}_{2}$ et $\mathrm{PGL}_{2}$ en
caractéristique 2. Une hypothèse est ajoutée au crayon quand l'exemple la
réclame. Et Grothendieck s'arrête, au milieu de la démonstration du théorème
principal, pour écrire que tel point est « à vérifier », au lieu de passer
outre. Une marge dit « rédiger mieux », une autre « !! ».

\keywords{SGA 3, smooth group scheme, group of multiplicative type,
maximal central torus, Rosenlicht unit lemma, character group, maximal
multiplicative quotient, extension by a torus, Picard group of an algebraic
group, simply connected semisimple group, Cartan subalgebra, regular element,
nilpotent Lie algebra, Engel's theorem, adjoint action, characteristic 2
counterexamples, representability of the centre, Chevalley structure theorem,
Cartan subgroup, abelian rank, reductive rank, unipotent rank, semisimple rank,
upper semicontinuity, Tate module, Frobenius weights, semi-abelian variety}
\end{resume}

\pagerange{1}{36}
\section*{Le fil du dossier, et les conventions}

\subsection*{Les stations}

Le dossier ne suit pas une démonstration, mais un traité : chaque station se
rattache à un exposé de SGA~3, et plusieurs se répondent.

\begin{itemize}
\item \textbf{Le tapuscrit de l'exposé XIII} (pages 2 à 6). Un groupe $G$
opère sur $V$, $W$ est un sous-schéma, et l'on compare $V$ au « balayage » de
$W$ par $G$ au moyen d'une application $\psi$. On le lit pour le cadre qu'il
fixe et pour ses corrections.
\item \textbf{Le tore central maximal} (pages 8 à 10) : plan d'un §6 pour
l'exposé XII, le lemme de Rosenlicht et ses corollaires, le quotient de type
multiplicatif universel.
\item \textbf{Prolonger des sous-groupes de type multiplicatif} (pages 11 et
12), au-dessus d'une base normale, et à partir de leurs points de torsion.
\item \textbf{Le groupe de Picard d'un groupe affine} (pages 13 à 15).
\item \textbf{Sous-algèbres de Cartan et éléments réguliers} (pages 17 à 19 et
21 à 26). C'est le cadre du tapuscrit appliqué à l'action adjointe, avec
$V = W(\mathfrak{g})$ et $W = W(\mathfrak{h})$. On y trouve onze conditions
sur une sous-algèbre $\mathfrak{h}$, un lemme sur les algèbres de Lie
nilpotentes, trois contre-exemples, et un centre non représentable.
\item \textbf{Le théorème de semi-continuité} (pages 27 à 36). Les invariants
numériques des fibres d'un schéma en groupes lisse varient
semi-continûment. La démonstration réutilise la technique des pages 11 et
12 : prolonger les points de torsion en sous-groupes finis étales.
\end{itemize}

Les pages 1 et 27 sont des chemises, dont l'inscription donne un titre. Les
pages 7 et 16 sont blanches. Les pages 20 et 30 sont des versos dactylographiés
où rien n'est de sa main : le premier traite de sous-algèbres nilpotentes
égales à leur normalisateur, comme la page 19 ; le second d'une lettre sur
les modules de différentielles.

\subsection*{Les conventions}

\emph{Groupes.} On écrit $\mathbf{G}_{m}$, $\mathbf{G}_{a}$ ; « de type
multiplicatif » en entier, là où il écrit « t.m. » ; $X^{*}(T)$ pour le groupe
des caractères, qu'il note $\mathbf{D}(T)$. Le manuscrit dit « préschéma » ;
on écrit « schéma ».

\emph{Le cadre du tapuscrit et de l'algèbre de Lie.} $N$ est le stabilisateur
de $W$ dans $G$, $M_{a}$ le transporteur de $a$ dans $W$, $\mathfrak{n}$ et
$\mathfrak{m}_{a}$ leurs algèbres de Lie. Pour une sous-algèbre $\mathfrak{h}$
de $\mathfrak{g}$, $\mathfrak{u}$ est son normalisateur dans $\mathfrak{g}$, et
$W(\mathfrak{g})$ le schéma vectoriel de $\mathfrak{g}$. Les mêmes lettres
servent dans les deux stations, et c'est voulu : la seconde est un cas
particulier de la première.

\emph{Les invariants.} Pour un groupe algébrique lisse et connexe $G$ sur un
corps algébriquement clos, on écrit $1 \to G_{\mathrm{aff}} \to G \to A \to 1$
sa décomposition de Chevalley, et l'on fixe, une fois pour toutes, les noms de
la page 29 plutôt que ceux, un peu différents, des pages 28 et 31. On en donne
la liste au début de la dernière section. Un invariant sans accent se rapporte
à la fibre spéciale, un invariant accentué ($\rho'$, $d'$) à la fibre
générique. $s_{0}$ est le point fermé, $s_{1}$ le point générique.

\emph{Deux homonymies à ne pas confondre.} La lettre $H$ désigne sur ces
pages trois objets différents. On garde $H(n)$ pour les familles de
sous-groupes finis des pages 12 et 33 à 36. On écrit $G^{\mathrm{tm}}$ le
quotient de type multiplicatif universel de la page 10, et $G_{\mathrm{aff}}$
la partie affine de la page 28. La lettre $R$ est un tore à la page 10 ; le
radical de la page 36 est noté $\mathrm{Rad}$.

\subsection*{Ce que le dossier annonce et n'établit pas}

\begin{itemize}
\item les §§ 6.1 à 6.3 de la page 8, qui ne sont qu'un plan, et le 6.4 annulé ;
\item la suite exacte de cochaînes de Hochschild de la page 10 ;
\item le corollaire 5.10 annulé (page 11) et la démonstration du corollaire de
la page 12 ;
\item les conditions c) et d) de la page 14, trop raturées pour être arrêtées ;
\item la liste de quatorze conditions à laquelle renvoient les pages 17 et 19,
qui n'est pas dans le dossier, et la remarque encadrée de la page 19 ;
\item le contre-exemple annulé de la page 26 ;
\item la représentabilité du centre annoncée à la page 25 ;
\item le point « à vérifier » de la page 33, dont dépend l'inégalité (5) ;
\item l'équivalence (3) $\Leftrightarrow$ (9) des pages 29 et 36, et la
réduction à une base de type fini sur $\mathbf{Z}$ (page 35).
\end{itemize}

\pagerange{2}{6}
\section*{Le tapuscrit de l'exposé XIII (pages 2 à 6)}

Cinq feuillets dactylographiés, paginés 6 à 10, portent au crayon « Exp XIII »
sur le premier.\footnote{Le premier mot de l'inscription, lu « Co », est
douteux. L'identification à l'exposé XIII de SGA~3 repose sur cette
inscription et sur la numérotation des énoncés. On ne confronte pas ces
feuillets au texte publié, et l'on ne dit pas si ses corrections y ont
passé.} Le texte tapé est publié. On en donne la teneur, sans le recomposer,
et l'on dit ce qu'il y a changé. Les lettres gothiques, les $\psi$ et
$\varphi$, et toutes les flèches d'implication sont de sa main, portées dans
les blancs de la frappe.

\emph{Le cadre.} $k$ est un corps, $G$ un $k$-groupe algébrique qui opère sur
un $k$-schéma $V$, et $W$ un sous-schéma de $V$. $N$ est le stabilisateur de
$W$, et $X = G \times^{N} W$ le produit contracté, avec
$p : X \to G/N$. Soit
\[
  \varphi : G \times W \to V, \quad (g, w) \mapsto g\,w,
\]
et $\psi : X \to V$ le morphisme qu'il induit. On note $e^{*}$ l'image de $e$
dans $G/N$. Pour $a \in W(k)$, on note $M_{a}$ le transporteur
$\{ g \in G : g\,a \in W \}$ ; alors $\varphi^{-1}(a) \simeq M_{a}$ et
$N \subset M_{a}$. On note $\mathfrak{n} \subset \mathfrak{m}_{a}$ leurs
algèbres de Lie.\footnote{Les définitions de $X$, $N$, $M_{a}$ sont sur les
pages 1 à 5 du tapuscrit, qui ne sont pas dans le dossier ; on les restitue.
C'est la seule lecture compatible avec ce que portent les feuillets :
l'isomorphisme $\varphi^{-1}(a) \simeq M_{a}$ (page 4), la projection plate
$p : X \to G/N$ (page 6), le point $(e^{*}, a)$ de $X$, et la condition
« $M_{a}$ et $N$ coïncident au voisinage de $e$ ». L'inégalité (2), qui
figure dans les hypothèses, est elle aussi sur une page absente.}

\textbf{Proposition 1.3.} On suppose $G$, $V$, $W$ géométriquement
irréductibles et l'inégalité (2) satisfaite. On considère les conditions
suivantes : (i) $\mathfrak{n} = \mathfrak{m}_{a}$ et $N$ lisse ;
(ii) $\mathfrak{n} = \mathfrak{m}_{a}$ et $\dim X = \dim V$ ;
(iii), (iii bis), (iii ter) $\mathfrak{n} = \mathfrak{m}_{a}$ et $\psi$
génériquement quasi-fini, resp.\ non ramifié, resp.\ étale ; (iv), (iv bis)
$\psi$ non ramifié, resp.\ étale, en $(e^{*}, a)$ ; (v) $\psi$ lisse en
$(e^{*}, a)$, $N$ lisse, et (2) une égalité ; (vi) $M_{a}$ et $N$ coïncident
au voisinage de $e$. On a les implications
\[
  \mathrm{(v)} \Leftrightarrow \mathrm{(iv\,bis)} \Rightarrow
  \mathrm{(iii\,ter)} \Rightarrow \mathrm{(i)} \Leftrightarrow
  \mathrm{(ii)} \Leftrightarrow \mathrm{(iii)} \Leftrightarrow
  \mathrm{(iii\,bis)} \Leftrightarrow \mathrm{(iv)} \Leftrightarrow
  \mathrm{(vi)} .
\]
Si $V$ est génériquement réduit, (iii ter) $\Leftrightarrow$ (i). Si de plus
$V$ est normal en $a$, toutes ces conditions sont équivalentes.

La chaîne d'implications est écrite de sa main, et c'est lui qui y insère
(iii bis). Il corrige surtout les hypothèses des deux dernières
phrases.\footnote{Le tapuscrit demandait « $V$ lisse sur $k$ et $W$
géométriquement réduit », puis « de plus $V$ géométriquement normal ». Il
écrit en interligne « génériquement » et « enfin … normal en $a$ ». La
démonstration de la page 3 n'utilise en effet que cela : $V$ génériquement
réduit pour passer de (iii bis) à (iii ter), car un morphisme non ramifié
au-dessus d'un corps est étale ; et $V$ normal en $a$ pour passer de (iv) à
(iv bis), par SGA~1, I~9.5. La correction ajuste l'énoncé à la preuve.}

\textbf{Corollaire 1.4.} Si $\mathfrak{n} = \mathfrak{m}_{a}$ et si $V$ est
normal en $a$, alors $\varphi$ est lisse en $(e, a)$. Par suite $M_{a}$ est
lisse sur $k$ en $e$, et $\psi$ est lisse en $(e^{*}, a)$. La démonstration
repose sur la chaîne
\[
  \dim V + \dim \mathfrak{n} \leqslant \dim (G \times W) \leqslant
  \dim_{(e,a)} \varphi^{-1}(a) + \dim V \leqslant \dim \mathfrak{m}_{a} +
  \dim V ,
\]
qui devient une suite d'égalités quand $\mathfrak{n} = \mathfrak{m}_{a}$. La
dernière est une égalité si et seulement si $M_{a}$ est régulier en $e$,
c'est-à-dire lisse sur $k$ en $e$.\footnote{« régulier en $e$ i.e. » est de sa
main, en interligne. L'équivalence est juste, parce que $e$ est un point
rationnel : en un point rationnel, régulier et lisse coïncident, et
$\dim_{e} M_{a} = \dim T_{e} M_{a}$ caractérise l'un et l'autre. Le tapuscrit
ajoute que la démonstration invoque des résultats délicats (la théorie de la
dimension, un théorème de Chevalley, EGA~IV pour la lissité) mais n'est pas
utilisée ensuite.}

\textbf{Corollaire 1.5.} C'est la version qui ne dépend pas du choix de $a$.
On suppose $V$ \emph{génériquement séparable sur $k$}.\footnote{Le tapuscrit
disait « génériquement réduit ». Il biffe et écrit « séparable sur $k$ ». La
correction est nécessaire : la démonstration commence par se ramener au cas
$k$ algébriquement clos, et l'hypothèse doit survivre à ce changement de base.
« Génériquement réduit » ne survit pas à une extension inséparable ; la
séparabilité du corps des fonctions (la réduction géométrique au point
générique) survit.} Les conditions (i), (i bis), (i ter) disent que $\psi$
est génériquement non ramifié, resp.\ étale, resp.\ lisse avec $N$ lisse et
(2) une égalité. Les conditions (ii) à (iv) disent qu'il existe $a \in W(k)$
satisfaisant la condition correspondante de 1.3. On a
(ii) $\Leftrightarrow$ (iii) $\Leftrightarrow$ (iv) $\Rightarrow$ (i)
$\Leftrightarrow$ (i bis) $\Leftrightarrow$ (i ter), et toutes sont
équivalentes si $k$ est algébriquement clos.\footnote{Dans (ii) et (iii), les
lettres gothiques portées à la main se lisent
« $\mathfrak{m} = \mathfrak{n}_{a}$ », dans l'ordre inverse de la page 2. On
lit $\mathfrak{n} = \mathfrak{m}_{a}$, seule forme qui ait un sens ici.}

\textbf{Corollaire 1.6.} Ici $V$ est normal et $k$ algébriquement
clos.\footnote{« $V$ normal, $k$ » est de sa main, devant l'interligne tapé
« alg.\ clos ».} Pour que $\psi$ soit \emph{birationnel}, il faut et il suffit
que les conditions de 1.5 soient satisfaites et qu'il existe un ouvert non vide
$U$ de $V$ tel que tout $a \in U(k)$ soit contenu dans un unique conjugué de
$W$. Pour $a \in V(k)$, les conditions suivantes sont alors équivalentes :
(i) un seul conjugué de $W$ contient $a$ ; (ii) les conjugués de $W$ qui
contiennent $a$ sont en nombre fini non nul ; (iii) $\psi^{-1}(a)$ a un point
isolé. Si de plus $W$ est réduit, et si $U$ est le plus grand ouvert au-dessus
duquel $\psi$ est un isomorphisme, elles équivalent aussi à (iv)
$a \in U(k)$. La démonstration passe par le \emph{Main Theorem} de Zariski et
la platitude de $p$.

\textbf{Remarque 1.7.} Sous les hypothèses de 1.6, un point satisfaisant (i)
à (iii) est dit \emph{$W$-régulier}. Sans supposer $k$ algébriquement clos, on
dit que $a \in V$ est $W$-régulier si $\psi^{-1}(a)$ est non vide et fini sur
$k(a)$.\footnote{La transcription, qui résume ce feuillet sans le recomposer,
porte ici $\varphi^{-1}(a)$. Or $\varphi^{-1}(a) \simeq M_{a}$ contient une
copie de $N$ (page 4), et ne peut être fini que si $N$ l'est. La définition
voulue, cohérente avec la condition (iii) de 1.6, porte sur $\psi^{-1}(a)$.
Nous ne savons pas si l'écart est dans le tapuscrit ou dans son résumé. « non
vide et » est tapé en interligne et relié à l'encre.} Sous les hypothèses plus
générales de 1.5, on introduit le degré $d$ de $\psi$ et le plus grand ouvert
au-dessus duquel $\psi$ est étale et fini. Il biffe la phrase selon laquelle
la notion serait « manifestement » stable par extension du corps de base à un
corps algébriquement clos.

\pagerange{8}{10}
\section*{Complément à l'exposé XII : le tore central maximal (pages 8 à 10)}

\emph{Le plan.} Un §6, « Le tore central maximal », puis des « tores quotients
d'un groupe algébrique ». 6.1 : définition en général, compatibilité avec
l'extension de la base, descente. 6.2 : existence lorsqu'il existe un
centre réductif. 6.3 : le cas d'un groupe algébrique.\footnote{Titres de sa
main, soulignés ; la fin de la ligne de titre et plusieurs mots de 6.2 et 6.3
sont illisibles, et une note « N.B. le groupe algébrique … affine » est en
partie perdue. Un 6.4 est rédigé puis annulé par quatre diagonales : pour $G$
lisse sur un corps parfait et $R$ son tore central maximal, il existerait un
sous-groupe fini $Z$ de $R$ tel que $G/Z$ se décompose à l'aide de $R/Z$ et
d'un sous-groupe $H$. La lecture des quotients et du mot qui dit comment ils
se combinent est douteuse ; on ne reconstruit pas l'énoncé.}

\textbf{Théorème 6.4 (Rosenlicht).} Soient $X$ et $Y$ deux $k$-schémas de
type fini géométriquement intègres, $a \in X(k)$ et $b \in Y(k)$. Toute
fonction inversible $u$ sur $X \times_{k} Y$ s'écrit
\[
  u(x, y) = u_{1}(x)\, u_{2}(y), \qquad u_{1}, u_{2} \text{ inversibles},
\]
et l'on peut prendre $u_{1}(x) = u(x, b)$, $u_{2}(y) = u(a, y)$ si l'on
normalise $u(a, b) = 1$.\footnote{La page écrit « $\varphi(x,y) = 1$ » ; la
lettre de l'argument est surchargée et ne se lit pas sûrement. La
normalisation en un point rationnel est la seule lecture qui donne un sens à
l'égalité. Le nom de Rosenlicht est sur la page, entouré au crayon ; le
résultat est de 1961.} Sous forme équivalente : si $u(x, b) = 1$ et
$u(a, y) = 1$ identiquement, alors $u = 1$. Sous cette forme on peut supposer
$k$ algébriquement clos, puis $X$ et $Y$ affines, et une réduction immédiate
ramène au cas où ils sont de dimension 1 : deux points d'une variété
irréductible sont toujours sur une courbe irréductible.\footnote{La page ne va
pas plus loin que cette réduction. Une ligne biffée,
« $X \times Y \subset \overline{X} \times \overline{Y}$ », indique la suite
attendue : plonger les courbes dans leurs complétées lisses et raisonner sur
les diviseurs de $u$. Tout un passage intermédiaire est très raturé.}

\textbf{Corollaire.} Soit $G$ un $k$-groupe lisse et connexe. Tout morphisme de
schémas $u : G \to \mathbf{G}_{m}$ tel que $u(e) = 1$ est un homomorphisme de
groupes. En effet, $G$ est géométriquement intègre, et
$(x, y) \mapsto u(xy)\, u(x)^{-1} u(y)^{-1}$ vaut $1$ sur $G \times \{e\}$ et
sur $\{e\} \times G$.\footnote{La démonstration est la nôtre ; la page énonce
le corollaire et souligne « lisse connexe ». Un corollaire précédent,
« Kif-kif, $G$ remplacé par un groupe … », est illisible.} En particulier le
groupe des caractères $X^{*}(G) = \mathrm{Hom}_{k\text{-gr}}(G, \mathbf{G}_{m})$
s'identifie au groupe des fonctions inversibles sur $G$ qui valent $1$ en
$e$. C'est un groupe abélien de type fini.

\textbf{Proposition.} Soit $G$ un $k$-groupe algébrique.
\begin{enumerate}
\item[a)] Il existe un épimorphisme $G \to G^{\mathrm{tm}}$ vers un groupe de
type multiplicatif, par lequel tout homomorphisme de $G$ vers un groupe de
type multiplicatif se factorise de façon unique.
\item[b)] Si $k$ est parfait, la formation de $G^{\mathrm{tm}}$ commute à
l'extension de la base.
\item[c)] Si $G$ est connexe \emph{et lisse}, $G^{\mathrm{tm}}$ est un tore.
\end{enumerate}
Si $k$ est imparfait, la formation de $G^{\mathrm{tm}}$ ne commute pas en
général à l'extension de la base, même pour $G$ connexe et
lisse.\footnote{« et lisse » est ajouté au crayon dans c), et il est
nécessaire : $G = \mu_{p}$ est connexe, et $G^{\mathrm{tm}} = \mu_{p}$ n'est
pas un tore. Le N.B.\ sur les corps imparfaits est juste, et en voici un
exemple, qui est le nôtre. Soit $k'= k(t^{1/p})$ avec $t \in k$ non puissance
$p$-ième, et $G = R_{k'/k}\mathbf{G}_{m}$ la restriction de Weil. Sur $k$,
$X^{*}(G)$ est engendré par la norme $N$, qui vaut $x \mapsto x^{p}$ sur le
tore diagonal $\mathbf{G}_{m} \subset G$. Sur $\bar k$, $G$ devient le produit
de $\mathbf{G}_{m}$ par un groupe unipotent, et $X^{*}(G_{\bar k})$ est
engendré par la projection $\chi$ sur $\mathbf{G}_{m}$, avec $N = \chi^{p}$.
Le quotient universel sur $k$, $N : G \to \mathbf{G}_{m}$, devient après
extension le composé de $\chi$ et du Frobenius, qui n'est pas le quotient
universel sur $\bar k$. Un diagramme en marge relie un tore $T \subset G$, un
sous-groupe $R \subset Z$ et un quotient de $G$ ; il est en partie
illisible.}

\textbf{Proposition.} Soient $G$ un $k$-groupe lisse et connexe et $R$ un
tore. Toute extension de $G$ par $R$ qui est triviale comme $R$-torseur est
triviale comme extension.

L'argument est celui de Rosenlicht. Le groupe connexe $G$ agit trivialement
sur le tore $R$, dont le groupe d'automorphismes est discret : l'extension
est centrale. Une section $s$ du torseur, normalisée par $s(e) = e$, donne un
cocycle $c(x, y) = s(x)\, s(y)\, s(xy)^{-1}$, morphisme de $G \times G$ dans
$R$, égal à $e$ sur $G \times \{e\}$ et $\{e\} \times G$. Il est donc
identiquement égal à $e$, et $s$ est un homomorphisme.\footnote{Énoncé et
démonstration sont reconstruits : la page porte « les extensions … de $G$ par
$R$ qui splittent … principaux homogènes splittent » et « En effet, d'ext.\
centrale … cocycles … d'une extension de $H$ par $R$, donc splitte », avec de
nombreux mots illisibles. « lisse » est ajouté par nous : le lemme de
Rosenlicht demande $G$ géométriquement intègre. Suivent, fermées par un
crochet, une suite exacte de complexes de cochaînes
$0 \to C^{\cdot}(e, R) \to C^{\cdot}(G^{\mathrm{tm}}, R) \to
\mathrm{Hom}_{\mathrm{gr}}(\ldots, R) \to 0$ et une cohomologie de
Hochschild biffée ; on ne les interprète pas.}

\pagerange{11}{12}
\section*{Prolonger un sous-groupe de type multiplicatif (pages 11 et 12)}

Deux corollaires pour l'exposé XII, puis un pour l'exposé XI. Le groupe $G$ y
est un $S$-schéma en groupes dont le feuillet ne rappelle pas les
hypothèses.\footnote{Dans les exposés XI et XII, où ces énoncés s'insèrent, $G$
est lisse et affine sur $S$ ; on le suppose. Cette lecture ne vérifie pas que
cela suffise à chacun des énoncés qui suivent.}

\textbf{Corollaire 5.8.} Soit $S$ noethérien, normal et irréductible, et $U$
un ouvert de $S$ dont le complémentaire est de codimension $\geqslant 2$. Tout
sous-groupe de type multiplicatif de $G_{U}$ provient d'un unique sous-groupe
de type multiplicatif de $G$.\footnote{Ce qui rend la normalité suffisante,
c'est que le prolongement a lieu \emph{dans} $G$. Pour un tore abstrait
au-dessus de $U$, l'énoncé analogue exige davantage. Voici un exemple, le
nôtre : soit $S$ l'hensélisé du cône $xy = z^{2}$ au sommet, sur un corps
algébriquement clos de caractéristique $\neq 2$. Il est normal, et
l'ouvert $U$ complémentaire du sommet a un revêtement étale connexe de degré
2. Le tore de normes associé n'a pas de prolongement en un tore sur $S$, car
$S$ est strictement hensélien et tout tore sur $S$ serait déployé.}

\textbf{Corollaire 5.9.} Soit $S$ normal, de point générique $y$, et
$\overline{H}_{y}$ un sous-tore de $G_{y}$. Pour que $\overline{H}_{y}$ soit la
fibre générique d'un sous-tore $\overline{H}$ de $G$, il faut et il suffit
que, pour tout $s \in S$ tel que $\dim \mathcal{O}_{S,s} = 1$, la composante
neutre réduite de la fibre en $s$ de l'adhérence de $\overline{H}_{y}$ soit de
type multiplicatif.\footnote{Il corrige « sous-groupe de t.m. » en
« sous-tore » et « un tore » en « de t.m. » ; pour un groupe réduit et
connexe, c'est la même chose. Une ou deux hypothèses sur $S$ sont
illisibles. En marge, de biais : « rédiger mieux ». L'énoncé combine un
critère aux points de codimension 1 et le corollaire 5.8 pour la codimension
$\geqslant 2$. Un corollaire 5.10, qui supposait $\overline{H}_{y}$ un tore,
est annulé.}

\textbf{Pour l'exposé XI, corollaire 4.?} Soient $S$ local et $G$ lisse et
affine sur $S$. Soit $(H(n))_{n}$ une famille de sous-groupes de type
multiplicatif de $G$, avec ${}_{m}H(n) = H(m)$ pour $m \mid n$, dont les fibres
spéciales sont les ${}_{n}H_{s}$ d'un sous-groupe de type multiplicatif $H_{s}$
de $G_{s}$. Il existe alors un unique sous-groupe de type multiplicatif $H$ de
$G$ tel que ${}_{n}H = H(n)$ pour tout $n$.\footnote{Les relations entre
$H(n)$, $H(m)$ et $H$ sont mal lisibles : le signe entre $H(n)$ et $H(m)$ ne
se lit pas, celui entre $H$ et $H(n)$ porte peut-être un indice, et la page
écrit $H(n)_{s} = H_{s}$. On les restitue d'après les pages 33 et 34, où la
même famille apparaît en clair : « $H(m) = {}_{m}H(n)$ si $m \mid n$ ». Un mot
avant « $S$ », lu « fini », pourrait porter une hypothèse sur $S$ ; on ne le
restitue pas et l'on ne démontre pas l'énoncé. En tête de la page : « Lemme
utilisé au corollaire nouveau ? » et, en marge, « !! ».} Une
\emph{remarque} envisage de n'utiliser que les $n$ dont les diviseurs
premiers sont dans un ensemble $P$ non vide de nombres premiers. Il faut
alors que la torsion du type de $H_{s}$ soit première à $P$. On y applique une
variante de la démonstration précédente et un argument de densité.

\pagerange{13}{15}
\section*{Le groupe de Picard d'un groupe affine (pages 13 à 15)}

\textbf{Proposition.} Soit $G$ un groupe algébrique affine sur un corps
algébriquement clos $k$.
\begin{enumerate}
\item[(a)] $\mathrm{Pic}(G)$ est fini.
\item[(b)] Si $G_{\mathrm{réd}}$ est résoluble, $\mathrm{Pic}(G) = 0$.
\item[(c)] Si $G$ est lisse et connexe, $B$ un sous-groupe de Borel et
$T \subset B$ un tore maximal,
\[
  \mathrm{Pic}(G) \simeq \mathrm{Pic}(G/B) / \mathrm{Im}\, X^{*}(T) ,
\]
où $X^{*}(T) \simeq X^{*}(B) \to \mathrm{Pic}(G/B)$ associe à un caractère le
fibré en droites qu'il définit sur $G/B$.
\item[(d)] Si $G$ est semi-simple, $\mathrm{Pic}(G) = 0$ si et seulement si
$G$ est simplement connexe.
\end{enumerate}

L'hypothèse sur $G_{\mathrm{réd}}$ plutôt que sur $G$ est à sa place : $G$
étant affine, $\mathrm{Pic}(G) = \mathrm{Pic}(G_{\mathrm{réd}})$, les
obstructions à relever un fibré en droites à travers un nilpotent étant des
groupes de cohomologie cohérente d'un schéma affine. Pour (b), un groupe
résoluble lisse et connexe est, comme variété, un produit de
$\mathbf{G}_{m}$ et de droites affines. La formule (c) vient de la suite exacte
$X^{*}(G) \to X^{*}(B) \to \mathrm{Pic}(G/B) \to \mathrm{Pic}(G) \to
\mathrm{Pic}(B) = 0$.\footnote{Ces justifications sont les nôtres ; la page
énonce. La marge porte deux exemples, sous un trait : $G/R = GP(n)$ et
$G = GL(n)$, qu'on lit $\mathrm{PGL}_{n}$ et $\mathrm{GL}_{n}$. On a
$\mathrm{Pic}(\mathrm{GL}_{n}) = 0$ et $\mathrm{Pic}(\mathrm{PGL}_{n}) =
\mathbf{Z}/n$. Elle porte aussi un diagramme qui compare les deux calculs de
(c) pour $G$ et $G/R$, $R$ le centre.}

\textbf{Proposition.} Soit $G$ un groupe algébrique affine sur un corps $k$
quelconque. Il existe un entier $n > 0$ tel que $n\, \mathrm{Pic}(G) = 0$.

\begin{enumerate}
\item[1°)] Soit $k'$ une extension finie de $k$ sur laquelle
$(G_{k'})_{\mathrm{réd}}$ est lisse, les points de $G/G^{\circ}$ sont
rationnels, le radical de $(G_{k'})^{\circ}_{\mathrm{réd}}$ est défini et
déployé, et le quotient semi-simple est déployé. L'entier $[k':k]$ annule le
noyau de $\mathrm{Pic}(G) \to \mathrm{Pic}(G_{k'})$. $\mathrm{Pic}(G)$ est
donc annulé par $[k':k]$ multiplié par l'exposant de
$\mathrm{Pic}(G_{k'})$.\footnote{La page dit qu'« on peut prendre pour $n$ le
degré de toute extension finie $k'$ » ayant les propriétés a) à d). Prise à la
lettre, l'affirmation est fausse : pour $G = \mathrm{PGL}_{2}$ sur $k$
algébriquement clos, $k' = k$ convient, et l'on aurait
$\mathrm{Pic}(\mathrm{PGL}_{2}) = 0$, alors qu'il vaut $\mathbf{Z}/2$. Ce
que le degré annule, par l'argument de norme
($N_{k'/k} \circ \mathrm{res} = [k':k]$), c'est le noyau de la restriction ;
on énonce cela. Les conditions c) et d) sont très raturées. Le « déployé » de
d) et les mots sur $N(T')$, $C(T')$ sont d'une lecture d'ensemble douteuse. Un
bloc annulé donnait deux cas particuliers : $G$ unipotent connexe, où une
puissance de la caractéristique suffit, et $G$ simple.}
\item[2°)] On suppose $G$ lisse et connexe et
$\mathrm{Hom}_{k_{s}\text{-gr}}(G_{k_{s}}, \mathbf{G}_{m}) = 0$, où $k_{s}$ est
une clôture séparable de $k$ et $\Gamma = \mathrm{Gal}(k_{s}/k)$. Alors
\[
  \mathrm{Pic}(G) \simeq \mathrm{Pic}(G_{k_{s}})^{\Gamma} .
\]
Par le corollaire au lemme de Rosenlicht (page 9), les fonctions
inversibles sur $G_{k_{s}}$ sont les constantes. La suite spectrale de
Hochschild--Serre donne alors
$0 \to H^{1}(\Gamma, k_{s}^{*}) \to \mathrm{Pic}(G) \to
\mathrm{Pic}(G_{k_{s}})^{\Gamma} \to \mathrm{Br}(k) \to \mathrm{Br}(G)$. Le
premier terme est nul par le théorème 90 de Hilbert, et la dernière flèche
est injective, scindée par le point rationnel $e$.\footnote{La page écrit
$\mathrm{Pic}(G) \simeq \mathrm{Pic}(G_{\tilde k})$, avec $\tilde k$ la
clôture « séparable », soulignée deux fois, ou « algébrique », écrite
au-dessus. Aucun des deux mots ne l'emporte clairement. Sans les invariants,
l'isomorphisme est faux. Pour une forme extérieure $G$ de
$\mathrm{PGL}_{3}$, déployée par une extension quadratique $L/k$, le groupe de
Galois agit par $-1$ sur $\mathrm{Pic}(G_{k_{s}}) = \mathbf{Z}/3$, et
$\mathrm{Pic}(G) = 0$. Avec la clôture algébrique, l'énoncé serait faux
d'une autre manière : les groupes unipotents non déployés sur un corps
imparfait ont un $\mathrm{Pic}$ non nul que $\bar k$ tue. Sa condition
entre crochets, que $G$ soit de « … unipotent nul », va dans ce sens. La
démonstration est la nôtre ; l'énoncé sous cette forme est un lemme de
J.-J.~Sansuc (1981), postérieur à ces notes.}
\item[3°)] Si $G$ est unipotent, $\mathrm{Pic}(G)$ est annulé par une
puissance de la caractéristique $p$ ; il le donne pour connu. Si $G$ est
réductif, lisse et connexe, et sans caractère sur $\bar k$, $\mathrm{Pic}(G)$
est fini, sous-groupe de $\mathrm{Pic}(G_{k_{s}}) = \mathrm{Pic}(G_{\bar
k})$.\footnote{La page : « Si $G$ est de … unipotent nul et
$\mathrm{Hom}_{\bar k\text{-gr}}(G_{\bar k}, \mathbf{G}_{m\bar k}) = 0$, alors
$\mathrm{Pic}(G) \simeq \mathrm{Pic}(G_{\bar k}) =$ \emph{groupe fini} ». On
lit « de radical unipotent nul », c'est-à-dire réductif, et l'on remplace
l'isomorphisme par l'inclusion que donne 2°). Sur $k_{s}$, $G$ est déployé, et
son $\mathrm{Pic}$ ne change plus par extension algébriquement close.}
\end{enumerate}

\pagerange{21}{22}
\section*{Sous-algèbres de Cartan et éléments réguliers (pages 17 à 26)}

\subsection*{Le cadre, et onze conditions (pages 21 et 22)}

C'est le cadre du tapuscrit, appliqué à l'action adjointe. $G$ est lisse sur
$k$, d'algèbre de Lie $\mathfrak{g}$, et $\mathfrak{h}$ est une sous-algèbre.
On prend $V = W(\mathfrak{g})$ et $W = W(\mathfrak{h})$. $N$ est le
normalisateur de $\mathfrak{h}$ dans $G$, $X = G \times^{N} W(\mathfrak{h})$ et
$\psi : X \to W(\mathfrak{g})$. Pour $a \in \mathfrak{h}$,
$\mathfrak{m}_{a} = \{ x \in \mathfrak{g} : [x, a] \in \mathfrak{h} \}$. Soit
enfin $\mathfrak{u}$ le normalisateur de $\mathfrak{h}$ dans $\mathfrak{g}$ ;
c'est l'algèbre de Lie du \emph{schéma} $N$, qui peut n'être pas lisse. On a
\[
  \mathfrak{h} \subset \mathfrak{u} \subset \mathfrak{m}_{a} .
\]
\emph{Un encadré de la page 21} dit : $\varphi$ est lisse en $(e, a)$ si et
seulement si $\mathrm{ad}(a)$ est injectif sur $\mathfrak{g}/\mathfrak{h}$,
c'est-à-dire $\mathfrak{m}_{a} = \mathfrak{h}$ ; et cela entraîne
$\mathfrak{h} = \mathfrak{u}$. En effet, la différentielle de $\varphi$ en
$(e, a)$ est $(x, b) \mapsto [x, a] + b$, de
$\mathfrak{g} \times \mathfrak{h}$ dans $\mathfrak{g}$. Elle est surjective si
et seulement si $\mathrm{ad}(a)$ est surjectif, donc bijectif, sur l'espace de
dimension finie $\mathfrak{g}/\mathfrak{h}$.\footnote{La justification est la
nôtre. L'identification $V = W(\mathfrak{g})$, $W = W(\mathfrak{h})$ l'est
aussi ; c'est ce que portent les notations des pages 21 et 22 ($\psi : X \to
W(\mathfrak{g})$, $N$, $\mathfrak{m}_{a}$, $\varphi$ lisse en $(e, a)$).}

La page 22 dresse une liste de conditions sur
$\mathfrak{h}$ :\footnote{Tout le corps de la page 22 est annulé par deux
longues diagonales, et se lit néanmoins. Les numéros (x) et (xi) récrivent
(xiii) et (xii), qui sont biffés : la liste de quatorze conditions à laquelle
renvoient les pages 17 et 19 en est un état antérieur, qui n'est pas dans le
dossier. En marge, des accolades réunissent (i)--(iii) et (ix)--(xi), et des
flèches doubles verticales relient (v) à (ix).}

\begin{enumerate}
\item[(i)] $\mathfrak{h}$ contient une sous-algèbre de Cartan ;
\item[(ii)] $\mathfrak{h} = \mathfrak{u}$ et $\mathfrak{h}$ contient des
éléments réguliers de $\mathfrak{g}$ ;
\item[(iii)] il existe $a \in \mathfrak{h}$ tel que $\mathrm{ad}(a)$ soit
injectif sur $\mathfrak{g}/\mathfrak{h}$, c'est-à-dire
$\mathfrak{m}_{a} = \mathfrak{h}$ ;
\item[(iv)] $\psi$ est dominant et $\mathfrak{h} = \mathfrak{u}$ ;
\item[(v)] $\psi$ est génériquement lisse ;
\item[(vi)] $X \times W(\mathfrak{h}) \to W(\mathfrak{g})$ est génériquement
lisse ;
\item[(vii)] $\psi$ est dominant ;
\item[(viii)] $\psi$ est dominant et quasi-fini ;
\item[(ix)] $\psi$ est génériquement quasi-fini et $\mathfrak{h} =
\mathfrak{u}$ ;
\item[(x)] $\psi$ est dominant, et $\mathfrak{h}$ est l'algèbre de Lie d'un
sous-groupe lisse de $G$ ;
\item[(xi)] $\psi$ est dominant, et $N$ est lisse.
\end{enumerate}

Au-dessous, il cherche la bonne définition. $\mathfrak{h}$ serait
« minimale » parmi les sous-algèbres qui sont égales à leur normalisateur et
satisfont au théorème de densité, ou qui sont définies par un sous-groupe
lisse, ou qui possèdent un $a$ comme en (iii). La liste s'arrête sur un e)
biffé.\footnote{Lecture d'ensemble douteuse. « Satisfaire au théorème de
densité » veut dire, pensons-nous, que les conjugués de $\mathfrak{h}$
remplissent un ouvert dense de $\mathfrak{g}$, c'est-à-dire (vii). Le mot
« densité » est répété et souligné, et lu avec doute.}

\pagerange{18}{18}
\subsection*{Un lemme sur les algèbres de Lie nilpotentes (page 18)}

\textbf{Lemme.} Soit $\mathfrak{h}$ une algèbre de Lie nilpotente sur un corps
\emph{infini} $k$, opérant sur un espace vectoriel $V$ de dimension finie par
$\rho$. Si $\rho(x)$ n'est injectif pour aucun $x \in \mathfrak{h}$, il existe
$v \in V$ non nul tel que $\rho(\mathfrak{h})\, v = 0$.

Sur $\bar k$, $V$ se décompose en sous-espaces primaires $V_{i}$ pour des
fonctions poids $\lambda_{i}$ sur $\mathfrak{h}$. L'hypothèse dit que pour
tout $x \in \mathfrak{h}$, l'un des $\lambda_{i}(x)$ est nul. Par le principe
du prolongement des identités, l'un des $\lambda_{i}$ est identiquement nul.
Alors $\mathfrak{h}$ opère sur $V_{i}$ par endomorphismes nilpotents, et le
théorème d'Engel donne un $v \neq 0$ tué par $\mathfrak{h}$. Comme le noyau
commun des $\rho(x)$ est défini sur $k$, il est non nul sur $k$.

Les lignes au crayon du bas de la page disent pourquoi le prolongement des
identités s'applique : si $n_{i} = \dim V_{i}$, le polynôme caractéristique
de $\rho(x)$ sur $V_{i}$ est $(t - \lambda_{i}(x))^{n_{i}}$. Donc
$\lambda_{i}(x)^{n_{i}} = \pm \det(\rho(x)|_{V_{i}})$ est une fonction
polynomiale de $x$. Le produit de ces déterminants s'annule sur
$\mathfrak{h}(k)$, donc est nul, $k$ étant infini.\footnote{L'argument est
celui de la page, complété par nous. Le complément compte : en
caractéristique $p$, les poids d'une algèbre de Lie nilpotente ne sont pas
nécessairement linéaires, et ce sont leurs puissances qui sont polynomiales.
La page écrit « $\exists\, v \in \mathfrak{h}$ » ; on lit $v \in V$.}

Appliqué à une sous-algèbre de Cartan $\mathfrak{c}$, nilpotente et égale à
son normalisateur, opérant sur $\mathfrak{g}/\mathfrak{c}$, le lemme donne un
$a \in \mathfrak{c}$ tel que $\mathrm{ad}(a)$ soit injectif sur
$\mathfrak{g}/\mathfrak{c}$. Sinon, un vecteur tué par $\mathfrak{c}$
fournirait un élément du normalisateur hors de $\mathfrak{c}$. Le même $a$
convient pour toute $\mathfrak{h} \supset \mathfrak{c}$, car
$\mathfrak{g}/\mathfrak{h}$ est un quotient de $\mathfrak{g}/\mathfrak{c}$. Sur
un corps infini, (i) entraîne donc (iii).\footnote{Cette application est la
nôtre ; c'est, pensons-nous, la raison d'être du lemme à cet endroit. La
chaîne « (xiv) $\Rightarrow$ (i) $\Rightarrow$ (iii) $\Rightarrow$ (v) » de la
page 19 est écrite dans l'ancienne numérotation, et nous ne savons pas si ses
(i) et (iii) sont ceux de la page 22.}

\pagerange{21}{23}
\subsection*{Deux contre-exemples en caractéristique 2 (pages 21 et 23)}

\emph{1) $\mathrm{SL}_{2}$.} Même si $\psi$ est birationnel et si
$\mathfrak{h}$ provient d'un sous-groupe lisse, la condition (ix) peut être
en défaut. Soit $k$ algébriquement clos de caractéristique 2 et
$G = \mathrm{SL}_{2}$, de base $X, Y, H$ avec $[X, Y] = H$ et $H$ central.
Alors $\mathfrak{g}$ est nilpotente, et sa seule sous-algèbre de Cartan est
$\mathfrak{g}$ elle-même. Soit $B$ un Borel et $\mathfrak{h} =
\mathrm{Lie}(B) = kH + kX$. Tout élément de $\mathfrak{g}$ a une droite
propre, donc est dans l'algèbre de Lie d'un Borel :
$\bigcup_{g} \mathrm{Ad}(g)\,\mathfrak{h} = \mathfrak{g}$, et $\psi$ est
dominant. $\mathfrak{h}$ est un idéal, donc $\mathfrak{u} = \mathfrak{g}
\neq \mathfrak{h}$ : (ix) est en défaut, (i) aussi. Comme $B \subset N \neq
G$, $\dim N = \dim B = \dim \mathfrak{h}$, et
$\dim X = \dim G$ : $\psi$ est génériquement quasi-fini. Il est de plus
birationnel.\footnote{La page le dit à la fin d'une phrase dont le début est
perdu (« … ici que $\psi$ est birationnel »). Une ligne biffée affirmait
d'abord « $\varphi$ lisse en tout point, $\psi$ étale en tout point », ce qui
est faux ici : $\mathfrak{m}_{a} \supset \mathfrak{u} = \mathfrak{g}$. La
birationalité est juste, mais elle a une raison subtile, qui est la nôtre.
Un élément générique $Z = \bigl(\begin{smallmatrix} a & b \\ c & a
\end{smallmatrix}\bigr)$ de $\mathfrak{g}$ a une seule valeur propre, double,
$a + \sqrt{bc}$, et une seule droite propre, $[\sqrt{b} : \sqrt{c}]$. Cette
droite n'est définie que sur une extension purement inséparable de degré 2 de
$k(\mathfrak{g})$, et $G \times^{B} W(\mathfrak{h}) \to W(\mathfrak{g})$
est radiciel de degré 2. Mais le normalisateur $N$ n'est pas $B$ : c'est un
schéma non lisse, d'algèbre de Lie $\mathfrak{u} = \mathfrak{g}$ de dimension
3. $G/N$ est le tordu de Frobenius de $G/B$, où la droite propre a pour
coordonnées $[b : c]$, rationnelles. La fibre générique de $\psi$ est un
point rationnel, et $\psi$ est birationnel. C'est le même phénomène qui met
(xi), « $N$ lisse », en défaut.}

\emph{2) $\mathrm{PGL}_{2}$.} Même si $\psi$ est dominant et
$\mathfrak{h} = \mathfrak{u}$, la condition (x) peut être en défaut. Soit
$G = \mathrm{PGL}_{2}$, toujours en caractéristique 2, de base
$X, Y, H$ avec
\[
  [X, Y] = 2H = 0, \qquad [H, X] = X, \qquad [H, Y] = Y .
\]
La sous-algèbre de Cartan contenant $H$ est $\mathfrak{t} = kH$. Soit
\[
  \mathfrak{h} = kH + k(\lambda X + \mu Y), \qquad \lambda \mu \neq 0 .
\]
C'est une sous-algèbre, égale à son normalisateur, qui contient la
sous-algèbre de Cartan $\mathfrak{t}$ : (i) est satisfaite. $\psi$ est
dominant. Mais $\mathfrak{h}$ n'est l'algèbre de Lie d'aucun sous-groupe lisse
de $G$. Un tel sous-groupe serait de dimension 2 et connexe, donc un Borel. Il
contiendrait alors le tore $T$ d'algèbre de Lie $\mathfrak{t}$, et les deux
Borels qui contiennent $T$ ont pour algèbres de Lie $kH + kX$ et
$kH + kY$.\footnote{La page est très raturée, et l'on en restitue
l'argument. Les vérifications de ce que $\mathfrak{h}$ est égale à son
normalisateur et de ce que $\psi$ est dominant sont les nôtres. Pour la
seconde : sur $\mathfrak{pgl}_{2}$ en caractéristique 2, la trace est bien
définie, et les orbites génériques sont les fibres de la trace. La trace de
$\gamma H + s(\lambda X + \mu Y)$ est $\gamma$, qui parcourt $k$. La page
écrit le groupe « $\mathrm{Pl}(2, k)$ » et le tableau de la page 17
« $GP(2, k)$ ».}

\pagerange{17}{19}
\subsection*{Le tableau des implications (pages 17 et 19)}

La page 17 dresse, dans l'ancienne numérotation à quatorze conditions, un
tableau d'implications et de non-implications. Il le rature une première fois
et le récrit dans une accolade :
(vii) $\not\Rightarrow$ (vi) si $\psi$ est birationnel, par
$\mathrm{SL}_{2}$ en caractéristique 2 ; (x) $\not\Rightarrow$ (v) ;
(i) $\not\Rightarrow$ (x), par $\mathrm{PGL}_{2}$ en caractéristique 2. Il
note aussi (xii) $=$ (x) $+$ (vii) et (xiii) $=$ (xii) $+$ (v). Les deux
exemples sont ceux qui viennent d'être lus.\footnote{Un « (v) $\not\Rightarrow$
(i) ??? » est encadré. En marge, un calcul de dimension pour
$G = \mathrm{PGL}_{2}$ et $\dim W = 1$ : $\dim X = 3 - 1 + 2 = 4$. On ne
traduit pas le tableau dans la numérotation de la page 22, faute de la liste
à laquelle il renvoie.}

La page 19, au crayon, prépare un troisième exemple par un calcul de
racines. Soit $\mathfrak{L} \subset \mathfrak{g}$, $\alpha > 0$ la seule racine
positive telle que $\bar\alpha = 0$, et $\beta > 0$ une racine telle que
$[X_{\alpha}, X_{\beta}] \neq 0$. Alors $\mathfrak{L} + k\,\mathfrak{g}_{-\beta}$
serait son propre normalisateur et satisferait au théorème de densité, mais
ne contiendrait « que des … algèbres de Cartan ».\footnote{Encadré, et de
lecture douteuse sur les trois mots décisifs ; le qualificatif des algèbres
de Cartan est illisible. On ne sait pas ce que désignent $\mathfrak{L}$ et la
barre. On ne reconstruit pas l'exemple. La page 20, verso dactylographié sans
rien de sa main, traite de sous-algèbres nilpotentes égales à leur
normalisateur.}

\pagerange{24}{24}
\subsection*{Un groupe résoluble sans élément régulier rationnel (page 24)}

\textbf{Exemple.} Il existe un groupe lisse, connexe et résoluble $G$ sur
$\mathbf{F}_{p}$ dont l'algèbre de Lie n'a aucun élément régulier rationnel.

Soit $T = \mathbf{G}_{m}^{2}$, d'algèbre de Lie $\mathfrak{t}$, avec
$\mathfrak{t}(\mathbf{F}_{p}) = \mathbf{F}_{p}^{2}$. Soit $(\alpha_{i})_{i \in
I}$ une famille finie de formes linéaires non nulles sur $\mathbf{F}_{p}^{2}$
telles que
\[
  \bigcup_{i} \mathrm{Ker}\, \alpha_{i} = \mathbf{F}_{p}^{2} ,
\]
par exemple une forme pour chacune des $p + 1$ droites du plan, ce qui assure
aussi $\bigcap_{i} \mathrm{Ker}\, \alpha_{i} = 0$. Soient
$\chi_{i} : T \to \mathbf{G}_{m}$ des caractères dont les différentielles sont
les $\alpha_{i}$, et
\[
  G = T \ltimes U, \qquad U = \mathbf{G}_{a}^{I},
\]
$T$ opérant sur le $i$-ème facteur par $\chi_{i}$. Un élément $x$ de
$\mathfrak{t}$ est régulier si et seulement si aucun $\alpha_{i}(x)$ n'est
nul. Sinon, $\mathrm{ad}(x)$ tue une droite de $\mathrm{Lie}(U)$, et le nilespace
de $x$ dépasse la dimension de $\mathfrak{t}$. Aucun élément de
$\mathfrak{t}(\mathbf{F}_{p})$ n'est donc régulier. Or tout tore maximal de $G$
défini sur $\mathbf{F}_{p}$ est conjugué de $T$ par un élément de
$U(\mathbf{F}_{p})$. Un élément régulier rationnel serait donc conjugué d'un
élément de $\mathfrak{t}(\mathbf{F}_{p})$.\footnote{« rationnel » est ajouté par
nous. Sur $\bar{\mathbf{F}}_{p}$, un élément générique de $\mathfrak{t}$ est
régulier, et c'est la finitude du corps qui fait l'exemple. Il montre que
l'hypothèse « $k$ infini », sous laquelle on construit les sous-algèbres de
Cartan à partir des éléments réguliers, ne peut pas être levée, même pour un
groupe résoluble. La page établit aussi, par un argument en partie illisible,
que le centre de $G$ est trivial. La prose de la page, en particulier
« caractères infinitésimaux » et la note de biais en marge, est d'une lecture
douteuse ; le second membre de l'égalité a été récrit plusieurs fois.}

\pagerange{25}{26}
\subsection*{Un centre non représentable (page 25)}

\textbf{Exemple.} Soit $S$ le spectre d'un anneau de valuation discrète $A$,
de corps résiduel $k$, de corps des fractions $K$. Il existe un $S$-schéma en
groupes séparé, étale et de type fini dont le foncteur centre n'est pas
représentable.

On part d'un groupe fini $E$ de centre trivial et d'un sous-groupe
commutatif non trivial $F$, par exemple engendré par un élément $s \neq
e$.\footnote{« fini » est ajouté par nous : la page dit « groupe
ensembliste », et « de type fini » l'exige. Par exemple $E = S_{3}$ et
$F$ engendré par une transposition.} Soit $G$ le sous-groupe ouvert du groupe
constant $E_{S}$ obtenu en retirant de la fibre spéciale le fermé
$E_{k} - F_{k}$. Ses fibres sont
\[
  G_{0} = F_{k}, \qquad G_{1} = E_{K} .
\]
Posons $A_{n} = A/\mathfrak{m}^{n+1}$. Sur $A_{n}$, $G$ est le groupe constant
commutatif $F$, dont le centre est $F$ tout entier. Sur le complété
$\hat A$, une section de $G$ est un élément de $F$, et elle est centrale
seulement si elle commute à $E$ sur la fibre générique ; le centre de $E$ étant
trivial, c'est $e$. Donc
\[
  Z(\hat A) = \{e\} \neq F = \varprojlim_{n} Z(A_{n}) .
\]
Pour tout schéma $Z'$, au contraire, $Z'(\hat A) = \varprojlim Z'(A_{n})$ : un
système compatible de points se factorise par un même ouvert affine, et
$\hat A = \varprojlim A_{n}$. Le centre n'est donc pas
représentable.\footnote{La dernière phrase est la nôtre ; la page dit que le
foncteur « ne commute pas » à la limite projective des $A_{n}$. Suit un
N.B.\ inachevé : « Si $G$ est plat de présentation finie sur $S$, à fibres
connexes, alors $Z$ est représentable …… ». L'énoncé n'est pas démontré, et
cette lecture ne le vérifie pas.}

La page 26 est barrée de trois traits : un « Contrexemple ! » sur les
sous-algèbres de Cartan contenant une sous-algèbre $\mathfrak{s}$ donnée, dont
les deux énoncés sont en grande partie illisibles, et qu'il abandonne.

\pagerange{27}{27}
\section*{Le théorème de semi-continuité (pages 27 à 36)}

La chemise de la page 27 porte « Théorème de ½ continuité ».\footnote{La
fraction ½ est écrite en petit, pour « semi- ».} Les notes sont à l'encre noire
sur papier bistre, et sans pagination de sa main. En revanche, les pages 31 à
36 portent ses numéros entourés (1) à (10), fixés à la page 29.

\pagerange{28}{29}
\subsection*{Les invariants (pages 28 et 29)}

Soit $G$ un groupe algébrique lisse et connexe sur un corps algébriquement
clos, et
\[
  1 \to G_{\mathrm{aff}} \to G \to A \to 1
\]
sa décomposition de Chevalley, $G_{\mathrm{aff}}$ affine lisse connexe, $A$
variété abélienne. Soient $T$ un tore maximal, $\mathrm{Rad}$ le radical de
$G_{\mathrm{aff}}$, $C = Z_{G}(T)$ le sous-groupe de Cartan et
$C_{\mathrm{aff}} = Z_{G_{\mathrm{aff}}}(T)$, qui est nilpotent, produit de $T$
et d'un groupe unipotent. On pose :

\begin{itemize}
\item $d = \dim G$, $\rho_{\mathrm{ab}} = \dim A$,
$d_{\mathrm{aff}} = \dim G_{\mathrm{aff}}$ ;
\item $\rho_{\mathrm{r}} = \dim T$, le rang réductif, et
$\rho_{\mathrm{rig}} = \rho_{\mathrm{ab}} + \rho_{\mathrm{r}}$ ;
\item $\rho_{\mathrm{ss}}$ le rang de $G_{\mathrm{aff}}/\mathrm{Rad}$, le rang
semi-simple ; $d_{\mathrm{ss}} = \dim G_{\mathrm{aff}}/\mathrm{Rad}$ et
$d_{\mathrm{rad}} = \dim \mathrm{Rad}$ ;
\item $\rho_{\mathrm{na}} = \dim C_{\mathrm{aff}}$,
$\rho_{\mathrm{u}} = \rho_{\mathrm{na}} - \rho_{\mathrm{r}}$, le rang
unipotent, et $\rho_{\mathrm{n}} = \dim C$, le rang nilpotent.
\end{itemize}

On a
\[
  \rho_{\mathrm{n}} = \rho_{\mathrm{na}} + \rho_{\mathrm{ab}} =
  \rho_{\mathrm{r}} + \rho_{\mathrm{u}} + \rho_{\mathrm{ab}}, \qquad
  d = \rho_{\mathrm{ab}} + d_{\mathrm{aff}} = \rho_{\mathrm{ab}} +
  d_{\mathrm{ss}} + d_{\mathrm{rad}} ,
\]
et $\rho_{\mathrm{r}}$, $\rho_{\mathrm{u}}$, $\rho_{\mathrm{ab}}$,
$\rho_{\mathrm{na}}$, $\rho_{\mathrm{n}}$ sont les mêmes pour $G$ et pour un
sous-groupe de Cartan.\footnote{Les noms varient d'une page à l'autre, et on
fixe ceux de la page 29. La page 28 écrit $\rho_{\mathrm{aff}}$ pour
$d_{\mathrm{aff}}$, $\rho_{\mathrm{s}}$ ou $\rho_{\mathrm{rs}}$ pour
$\rho_{\mathrm{ss}}$, $d_{\mathrm{s}}$ pour $d_{\mathrm{ss}}$, et
$H$ pour $G_{\mathrm{aff}}$. Elle définit $\rho_{\mathrm{rad}}$ comme le rang
du radical, mais l'emploie deux lignes plus bas comme sa dimension, dans
« $\rho_{\mathrm{aff}} = \rho_{\mathrm{ss}} + \rho_{\mathrm{rad}}$ » ; on
écrit $d_{\mathrm{ss}}$ et $d_{\mathrm{rad}}$. Elle exprime aussi
$\rho_{\mathrm{u}}$ et $\rho_{\mathrm{na}}$ à l'aide d'un $V$ qu'elle ne
définit pas ; on le lit $C_{\mathrm{aff}}$. Que $C$ ait les mêmes invariants
que $G$ se voit ainsi : $C$ contient $T$, se surjecte sur $A$ (car
$G = G_{\mathrm{aff}} \cdot N_{G}(T)$), et est son propre sous-groupe de
Cartan. Le premier indice, « ab », se lit aussi « as ».}

La page 29 est une feuille de calculs. On y lit la liste numérotée qui servira
de repère :
\[
  \begin{array}{llll}
    (1)\ \rho_{\mathrm{ab}} & (2)\ \rho_{\mathrm{rig}} &
      (3)\ \rho_{\mathrm{ss}} & (4)\ d_{\mathrm{ss}} \\
    (5)\ \rho_{\mathrm{n}} & (6)\ \rho_{\mathrm{u}} &
      (7)\ \rho_{\mathrm{na}} & (8)\ d_{\mathrm{aff}}
  \end{array}
\]
puis (9) $d_{\mathrm{rad}}$ et (10) $d_{\mathrm{radaff}}$, ce dernier
non défini. Elle note les implications (2) $+$ (5) $\Rightarrow$ (6),
(1) $+$ (5) $\Rightarrow$ (7), (1) $\Leftrightarrow$ (8), (3) $\Leftrightarrow$
(9), (1) $+$ (3) $\Rightarrow$ (10). Les trois premières sont de
l'arithmétique : $d$ étant constant dans une famille lisse, on a
$\rho_{\mathrm{u}} = \rho_{\mathrm{n}} - \rho_{\mathrm{rig}}$,
$\rho_{\mathrm{na}} = \rho_{\mathrm{n}} - \rho_{\mathrm{ab}}$ et
$d_{\mathrm{aff}} = d - \rho_{\mathrm{ab}}$.\footnote{L'équivalence
(3) $\Leftrightarrow$ (9), qui revient à la page 36, n'est pas expliquée, et
cette lecture ne sait pas la justifier : la dimension du radical n'est pas
déterminée par le rang semi-simple. Le haut de la page porte des calculs
épars : ${}_{n}A_{0}$, $T_{0} \subset C_{0}$, « V.A. »,
$H^{2}(\ , T_{0}) \simeq H^{2}(\ , \mathbf{Q}/\mathbf{Z})$,
$\rho_{\mathrm{ab}} \leqslant \rho'_{\mathrm{ab}}$. Ils annoncent la
démonstration. Le premier chiffre de « (2) $+$ (5) » est surchargé.}

\pagerange{31}{32}
\subsection*{L'énoncé et les réductions (pages 31 et 32)}

\textbf{Théorème.} Soit $G$ un schéma en groupes lisse, séparé et de
présentation finie sur $S$. Pour $s \in S$, on évalue chaque invariant sur la
composante neutre de la fibre géométrique en $s$.
\begin{enumerate}
\item[(i)] $d_{\mathrm{aff}}$, $d$, $\rho_{\mathrm{n}}$, $\rho_{\mathrm{na}}$ et
$\rho_{\mathrm{u}}$ sont semi-continues supérieurement.
\item[(ii)] $\rho_{\mathrm{ab}}$, $\rho_{\mathrm{rig}} = \rho_{\mathrm{ab}} +
\rho_{\mathrm{r}}$ et $\rho_{\mathrm{ss}}$ sont semi-continues
inférieurement.
\end{enumerate}

Deux exemples font voir le sens des inégalités et pourquoi c'est
$\rho_{\mathrm{ab}} + \rho_{\mathrm{r}}$, et non $\rho_{\mathrm{r}}$, qui est
semi-continue. Une courbe elliptique à réduction multiplicative dégénère en
$\mathbf{G}_{m}$ : $\rho_{\mathrm{ab}}$ tombe de 1 à 0 et $\rho_{\mathrm{r}}$
monte de 0 à 1, leur somme reste 1. Le schéma en groupes
$\mathrm{Spec}\, A[x, (1 + \pi x)^{-1}]$ sur un anneau de valuation discrète
d'uniformisante $\pi$ a pour fibre générique $\mathbf{G}_{m}$ et pour fibre
spéciale $\mathbf{G}_{a}$ : $\rho_{\mathrm{r}}$ tombe, $\rho_{\mathrm{u}}$
monte.\footnote{Les deux exemples sont les nôtres. Dans (ii), « Si $G$ est
plat sur $S$ » est biffé, et « lisse » est ajouté en interligne à l'énoncé.
Les notations $\rho_{\mathrm{r}}$, $\rho_{\mathrm{u}}$, $\rho_{\mathrm{n}}$
sont, à notre connaissance, celles de l'exposé XII de SGA~3, qui traite le cas
affine. Ce qui est propre à ces pages est le cas général, avec la partie
abélienne, où $\rho_{\mathrm{r}}$ seul n'est plus semi-continu.}

\emph{Réductions.} « Comme d'habitude », on se ramène au cas où $S$ est le
spectre d'un anneau de valuation discrète, de point fermé $s_{0}$ et de point
générique $s_{1}$.\footnote{Ce « comme d'habitude » suppose que les fonctions
en jeu sont localement constructibles, ce qui permet de ramener la
semi-continuité à la spécialisation le long des traits. La page ne le dit
pas, et cette lecture ne le vérifie pas.} Il faut prouver, les invariants
accentués étant ceux de la fibre générique :
\[
  d_{\mathrm{aff}} \geqslant d'_{\mathrm{aff}}, \quad d \geqslant d', \quad
  \mathrm{(5)}\ \rho_{\mathrm{n}} \geqslant \rho'_{\mathrm{n}}, \quad
  \mathrm{(6)}\ \rho_{\mathrm{u}} \geqslant \rho'_{\mathrm{u}}, \quad
  \mathrm{(7)}\ \rho_{\mathrm{na}} \geqslant \rho'_{\mathrm{na}},
\]
\[
  \mathrm{(1)}\ \rho_{\mathrm{ab}} \leqslant \rho'_{\mathrm{ab}}, \quad
  \mathrm{(2)}\ \rho_{\mathrm{rig}} \leqslant \rho'_{\mathrm{rig}}, \quad
  \mathrm{(3)}\ \rho_{\mathrm{ss}} \leqslant \rho'_{\mathrm{ss}} .
\]
Ces quantités ne changent pas quand on remplace un groupe algébrique par sa
composante neutre réduite. De plus, $d_{\mathrm{aff}}$ et $d$ croissent avec
le groupe. On peut donc remplacer $G$ par l'adhérence de sa fibre générique,
qui est plate. Quitte à faire une extension finie de $S$, on peut supposer
$G_{1}$ connexe et lisse. On peut enfin supposer $G_{0}$ connexe, mais
peut-être pas lisse.\footnote{Les crochets de la page sont déséquilibrés : le
premier s'ouvre en bas de la page 31 et se ferme à la page 32. Plusieurs mots
de la phrase de réduction sont illisibles ou douteux (« quitte », « réserve
de platitude »).}

\textbf{a)} Pour $G$ lisse, $d = d'$ ; alors $d_{\mathrm{aff}} \geqslant
d'_{\mathrm{aff}}$ équivaut à $\rho_{\mathrm{ab}} \leqslant
\rho'_{\mathrm{ab}}$, qu'on prouve plus bas.

\pagerange{32}{33}
\subsection*{Le sous-groupe de Cartan : l'inégalité (5) (pages 32 et 33)}

\textbf{b)} On peut supposer que $G_{0}$ a un tore maximal déployé $T_{0}$. Soit
$C_{0}$ le sous-groupe de Cartan qu'il définit. Pour $n$ premier à la
caractéristique résiduelle, $S$ étant hensélien, le sous-groupe
${}_{n}T_{0}$ se prolonge en un sous-groupe fini étale $H(n)$ de $G$, avec
$H(m) = {}_{m}H(n)$ si $m \mid n$.\footnote{Une première voie est abandonnée
(lignes biffées et traversées de courbes). Elle passait par $C_{0}/T_{0}$,
extension d'une variété abélienne $A_{0}$ par un groupe unipotent lisse
connexe $U_{0}$, et par le scindage canonique des extensions de
${}_{n}A_{0}$ par $U_{0}$. En marge : « $G$ lisse$/S$ », puis « $G$ lisse$/k$ ».}
Les centralisateurs $\underline{\mathrm{Cent}}_{G}\, H(n)$ sont des
sous-groupes fermés lisses de $G$. Leur famille est stationnaire, et sa
valeur limite $C$ a pour fibre spéciale
\[
  C_{0} = \underline{\mathrm{Cent}}_{G_{0}}(T_{0}) ,
\]
la réunion des ${}_{n}T_{0}$ étant dense dans $T_{0}$.

Il faudrait aussi que $C_{\bar s_{1}}$ contienne un sous-groupe de Cartan de
$G_{\bar s_{1}}$. La page l'affirme et ajoute : « c'est un [point] à
vérifier ».\footnote{Souligné d'une ligne ondulée. C'est de cela que dépend
l'inégalité (5). Nous ne le vérifions pas non plus. La réunion des
$H(n)_{1}$ a pour adhérence un groupe commutatif dont la composante neutre
peut avoir une partie abélienne. Que son centralisateur contienne un
sous-groupe de Cartan n'est pas immédiat.} Admettons-le. Les invariants
$\rho_{\mathrm{r}}$, $\rho_{\mathrm{u}}$, $\rho_{\mathrm{ab}}$,
$\rho_{\mathrm{n}}$, $\rho_{\mathrm{na}}$ sont alors les mêmes pour $G_{1}$ et
$C_{1}$, et pour $G_{0}$ et $C_{0}$. Pour (5), on est donc ramené au cas
$G = C$, c'est-à-dire $G_{0}$ nilpotent. L'inégalité
\[
  \rho_{\mathrm{n}} \geqslant \rho'_{\mathrm{n}}
\]
devient triviale, puisque $\rho_{\mathrm{n}} = d = d' \geqslant
\rho'_{\mathrm{n}}$.\footnote{Le chiffre devant « formule » se lit « 1) » ou
« 5) ». La relation est encadrée sur la page.}

\pagerange{33}{34}
\subsection*{Un lemme sur les groupes nilpotents, et le passage au cas commutatif (pages 33 et 34)}

\textbf{Lemme.} Soit $C$ un groupe lisse, connexe et nilpotent sur un corps
algébriquement clos $k$ de caractéristique $p \geqslant 0$. Soit
$Z = \underline{\mathrm{Cent}}(C)^{\circ}_{\mathrm{réd}}$. Alors $C/Z$ est
unipotent. Pour tout groupe $H$ de type multiplicatif,
\[
  \underline{\mathrm{Hom}}_{k\text{-gr}}(H, C) \simeq
  \underline{\mathrm{Hom}}_{k\text{-gr}}(H, Z) .
\]
En particulier ${}_{n}C = {}_{n}Z$ pour $n$ premier à $p$.\footnote{La page
place $Z$ au-dessus d'une accolade qui semble couvrir tout le quotient
$C/\underline{\mathrm{Cent}}(C)^{\circ}_{\mathrm{réd}}$. Lu ainsi, le second
énoncé serait faux : il n'y a pas d'homomorphisme non trivial d'un groupe de
type multiplicatif vers un groupe unipotent, alors que $C$ contient son tore
maximal. On lit donc $Z$ comme le centre, conformément à ce que la suite en
fait. « premier à $p$ » est ajouté par nous. La justification est la nôtre :
$C$ est engendré par sa partie anti-affine, centrale, et par sa partie
affine, produit d'un tore central et d'un groupe unipotent ; une image de
type multiplicatif tombe dans le tore, et un point de torsion d'ordre
premier à $p$ a une image triviale dans le quotient unipotent.}

On l'applique à $C = G_{0}$. Si $S$ est local strictement hensélien, les
${}_{n}G_{0}$ se prolongent de façon unique en des sous-groupes finis étales
$H(n)$ de $G$, centraux, avec ${}_{m}H(n) = H(m)$ si $m \mid n$. Or
$\rho_{\mathrm{ab}}$, $\rho_{\mathrm{r}}$ et donc $\rho_{\mathrm{rig}}$
croissent avec le groupe. Soit $K$ l'adhérence dans $G$ de la réunion des
$H(n)_{1}$. On a
\[
  \rho_{\mathrm{ab}}(K_{1}) \leqslant \rho'_{\mathrm{ab}}, \qquad
  \rho_{\mathrm{rig}}(K_{1}) \leqslant \rho'_{\mathrm{rig}} ,
\]
et $\rho_{\mathrm{ab}}(K_{0}) = \rho_{\mathrm{ab}}(G_{0})$,
$\rho_{\mathrm{r}}(K_{0}) = \rho_{\mathrm{r}}(G_{0})$, puisque $K_{0}$ contient
tous les ${}_{n}G_{0}$. Pour (1) et (2), on peut donc remplacer $G$ par
$K$.\footnote{« [qu'on peut supposer lisse$/k(s_{1})$ !] » est ajouté en
interligne. La justification de l'égalité des invariants sur la fibre spéciale
est lue avec doute (« Cartan … »). Un « Corollaire » est biffé en tête de page.}

On a perdu la limite de $G$, mais gagné la \emph{commutativité}. Pour un groupe
commutatif, $\rho_{\mathrm{n}} = d$, donc
$\rho_{\mathrm{u}} = d - \rho_{\mathrm{rig}}$ et
$\rho_{\mathrm{na}} = d_{\mathrm{aff}} = d - \rho_{\mathrm{ab}}$. L'inégalité
(2) équivaut alors à (6), et (1) à (7) et à (8).

\pagerange{35}{36}
\subsection*{Le cas commutatif : unipotence et affinité (pages 35 et 36)}

Il reste à prouver, pour $G$ commutatif, $\rho_{\mathrm{u}} \geqslant
\rho'_{\mathrm{u}}$ et $d_{\mathrm{aff}} \geqslant d'_{\mathrm{aff}}$.

\emph{Unipotence.} En remplaçant $G$ par l'adhérence de la partie unipotente
de $G_{1}$, on est ramené à prouver que si $G_{1}$ est unipotent, $G_{0}$
l'est aussi. C'est « facile par les points d'ordre fini ». Si $G_{0}$ n'était
pas unipotent, il contiendrait un tore ou une partie abélienne, donc des
points d'ordre $n$ premier à la caractéristique résiduelle. Ceux-ci se
prolongent en sous-groupes finis étales, dont les fibres génériques seraient
des points d'ordre $n$ d'un groupe unipotent. Or il n'y en a pas.\footnote{Le
développement de « facile » est le nôtre. En marge, deux fois :
« [sans restriction … limitation] ».}

\emph{Affinité.} Pour $d_{\mathrm{aff}}$, on est ramené à prouver que si $G_{1}$
est affine, $G_{0}$ l'est aussi, c'est-à-dire $\rho_{\mathrm{ab}}(G_{0}) = 0$.
On se ramène au cas où $S$ est semi-local de dimension 1, localisé d'un schéma
intègre de type fini sur $\mathbf{Z}$. Le corps résiduel $k(s_{0})$ est alors
fini, et l'on peut supposer déployé le tore maximal de $G_{1}$.\footnote{La
réduction à une base de type fini sur $\mathbf{Z}$ est annoncée
(« à priori … ») et n'est pas conduite ; nous ne la vérifions pas. En marge,
un croquis : une flèche vers « $S$ », un angle droit, un segment pointé.}
Passons au hensélisé $\tilde S$ en $s_{0}$ — « mais pas strictement
hensélien ! » — et fixons un nombre premier $\ell$ distinct de la
caractéristique résiduelle. Pour $n = \ell^{m}$, les ${}_{n}G_{0}$ se
prolongent en sous-groupes finis étales $H(n)$ sur $\tilde S$. Le groupe
\[
  \pi_{1}(\tilde S) = \pi_{1}(s_{0}) \simeq \hat{\mathbf{Z}}
\]
opère sur $M_{n} = H(n)_{\bar s_{0}}$, donc sur
\[
  M = \varprojlim_{n} M_{n} \simeq
  \mathbf{Z}_{\ell}^{\,2\rho_{\mathrm{ab}} + \rho_{\mathrm{r}}} ,
\]
le module de Tate de la partie semi-abélienne de $G_{0}$.\footnote{La page
écrit $\pi_{1}(s_{0}) \simeq \mathbf{Z}$ ; c'est le groupe de Galois d'un corps
fini, $\hat{\mathbf{Z}}$, engendré topologiquement par le Frobenius.} D'autre
part, $M$ s'identifie, de façon compatible à l'action, à
$\varprojlim H(n)_{\bar s_{1}}$. Ce module est contenu dans
$\varprojlim ({}_{n}G)_{\bar s_{1}}$, qui, $G_{1}$ étant affine et
commutatif, est le module de Tate de son tore déployé :
$(\mathbf{Z}_{\ell}(1))^{r}$. Le groupe de Galois opère donc sur $M$ par le
\emph{caractère cyclotomique} : le Frobenius de $k(s_{0})$, de cardinal $q$,
y opère par multiplication par $q$.

Or si $\rho_{\mathrm{ab}} \neq 0$, $M$ a pour quotient le module de Tate de
la variété abélienne $A_{0}$ sur le corps fini $k(s_{0})$. Par le théorème de
Weil, le Frobenius y a des valeurs propres de valeur absolue $q^{1/2}$, et
non $q$. C'est impossible : $\rho_{\mathrm{ab}}(G_{0}) = 0$, et $G_{0}$ est
affine.\footnote{La fin de l'argument est reconstruite. La page dit que les
opérations de Galois se font « par l'intermédiaire exclusif » d'un
« caractère » (lecture douteuse), puis qu'on est ramené à un énoncé « sur une
V.A.\ … sur un corps $\mathfrak{k}$ de type fini … mais ici suffisant », avec
un crochet ouvrant absent. Nous comprenons que l'énoncé général, sur un
corps de type fini, n'est pas nécessaire, le corps résiduel étant ici fini. Le
recours au théorème de Weil sur les poids du Frobenius (1948) est le nôtre. Un
argument voisin, par le module de Tate de la partie semi-abélienne et
l'action de l'inertie, est celui de SGA~7, exposé IX (1967-1969) ; on ne sait
pas lequel des deux textes précède l'autre.}

\pagerange{36}{36}
\subsection*{Le rang semi-simple (page 36)}

« Restent à prouver (3) et (9) », avec la note « (3) $\Leftrightarrow$ (9) »
et un résultat sur la spécialisation des groupes réductifs.\footnote{La page
écrit que « un groupe réductif … spécialise en un réductif », un mot illisible
au milieu. Sous cette forme, c'est faux : le groupe
$\mathrm{Spec}\, A[x, (1 + \pi x)^{-1}]$ de la page 31 a pour fibre générique
$\mathbf{G}_{m}$ et pour fibre spéciale $\mathbf{G}_{a}$. Ce qui est vrai, pour
$G$ lisse et affine, c'est l'inverse : si la fibre spéciale est réductive, la
fibre générique l'est aussi. Le lieu des fibres réductives est ouvert (SGA~3,
exposé XIX, à notre connaissance). Le mot illisible est peut-être une
négation. On ne tranche pas.} Il esquisse la preuve de (3).

Soit $T_{0}$ un tore maximal de $G_{0}$, et $\mathrm{Rad}_{0}$, $\mathrm{Rad}_{1}$
les radicaux des parties affines des deux fibres. Soit $r_{0} =
\rho_{\mathrm{ss}}(G_{0})$ ; on a $r_{0} = \dim T_{0}/(T_{0} \cap
\mathrm{Rad}_{0})$. On prolonge les ${}_{n}T_{0}$ en une famille compatible de
sous-groupes $H(n)$ de $G$. L'adhérence de $\mathrm{Rad}_{1}$ est un sous-groupe
distingué résoluble ; sa fibre spéciale est donc contenue dans
$\mathrm{Rad}_{0}$. Les images des $H(n)_{1}$ dans
$G_{1}/\mathrm{Rad}_{1}$ forment alors un système compatible dont le rang,
quand $n = \ell^{m}$ croît, est au moins $\dim T_{0}/(T_{0} \cap
\mathrm{Rad}_{0}) = r_{0}$. Ce rang est au plus le rang de
$G_{1}/\mathrm{Rad}_{1}$, c'est-à-dire $\rho'_{\mathrm{ss}}$.\footnote{Le paragraphe est numéroté « (4) », et la plupart de ses mots de
liaison sont illisibles. L'argument qui borne la fibre spéciale de
l'adhérence de $\mathrm{Rad}_{1}$ est le nôtre, suggéré par la marge
« $\overline{R}_{1} = R$ ». La page parle de « $n$-rang ». Pour un seul $n$,
la borne par le rang est fausse : $(\mathbf{Z}/2)^{2}$ se plonge dans
$\mathrm{PGL}_{2}$, de rang 1, en caractéristique $\neq 2$. Elle vaut pour le
système compatible tout entier, dont l'adhérence est un groupe commutatif de
tore de dimension au moins le rang $\ell$-adique du système. C'est ainsi qu'on
lit.} D'où
\[
  \rho_{\mathrm{ss}} \leqslant \rho'_{\mathrm{ss}} .
\]

\end{document}
